\section{The characterwise elliptic-unit determinant and its numerical gate}
\label{sec:elliptic-unit-specialization}

The fixed-base reduction AC-K25 permits use of the
imaginary-quadratic main conjecture character by
character.  Its consequence here is an exact
determinant specialization, with the torsion term
retained.  Equality of determinants is not itself
an upper bound for their specialization indices.

\paragraph{Primary input and conventions.}
Johnson-Leung--Kings, \emph{On the equivariant and the
non-equivariant main conjecture for imaginary quadratic
fields}, arXiv:0804.2828v2, Theorem~5.2, p.~24, applies
characterwise at every prime.  Exact supporting locators
are Corollary~5.3, pp.~24--25; Lemmas~4.3, 4.4 and 4.7,
pp.~20--22; equation~(17), p.~24; and Lemma~6.6, p.~32.
The complete printed proof is \S6, pp.~29--33.
Its scalar inversion after equation~(14) is transported
below by $\gamma\mapsto\gamma^{-1}$.
The equivariant Conjecture~5.5 is not assumed.
\url{https://arxiv.org/pdf/0804.2828}.

\begin{lemma}[AC-K26: the exact specialized determinant index]
\label{AC-K26}
Let $O$ be a complete discrete valuation ring of
characteristic zero, with fraction field $E$, uniformizer
$\pi$, and finite residue field $k_O$.
Let $\Lambda=O[[T_1,T_2]]$, and let $\mathscr C$ be a
perfect complex of $\Lambda$-modules, with a specified
generator $z_\Lambda$ of its inverse determinant line
$\det_\Lambda^{-1}\mathscr C$.
For a continuous $O$-algebra specialization
$\alpha:\Lambda\to O$, put
$\mathscr C_\alpha=O\otimes^\mathbf L_{\Lambda,\alpha}\mathscr C$.
Suppose its only cohomology is in degrees $1,2$,
$H^1$ has rank one, and $H^2$ is finite.
Let $z_\alpha$ be the image of $z_\Lambda$ in
\[
 \det_O^{-1}\mathscr C_\alpha\otimes_O E
                       \simeq H^1(\mathscr C_\alpha)\otimes_O E.
\]
For a generator $u$ of $H^1(\mathscr C_\alpha)_{\rm free}$,
write $z_\alpha=d_\alpha u$, with $d_\alpha\in E^\times$,
and put $i_\alpha=v_\pi(d_\alpha)$.  Then
\[
 \operatorname{length}_O H^2(\mathscr C_\alpha)
   =i_\alpha+
           \operatorname{length}_O H^1(\mathscr C_\alpha)_{\rm tors}.
                                                            \tag{EU.1}
\]
The index $i_\alpha$ is an integer, independent of
the chosen generator $u$, and may be negative.
Its positive part obeys
\[
 0\leq \operatorname{length}_O H^2-\max(0,i_\alpha)
        \leq\operatorname{length}_O H^1_{\rm tors}.  \tag{EU.2}
\]
\end{lemma}
\begin{proof}
Represent $\mathscr C$ by a bounded complex of finite
free modules.  Formation of the alternating tensor
product of their top exterior powers commutes with
base change.  Thus $z_\Lambda$ specializes to a
generator of the inverse determinant line over $O$;
in particular its rational image is nonzero.
This uses the derived complex throughout, so no
cohomological Tor term is discarded.

Over a discrete valuation ring, each finite module
is a direct sum of $O/\pi^aO$.
The resolution $O\xrightarrow{\pi^a}O$ shows that
a length-$a$ torsion group in degree $2$ contributes
$\pi^aO$ to the inverse determinant lattice, whereas
one in degree $1$ contributes $\pi^{-a}O$.
These signs can also be checked by trivializing the
two-term complex over $E$: in degrees $1,2$ its
inverse determinant is multiplication by $\pi^a$.
The determinant is multiplicative for successive
cohomology truncations, as follows by cancellation
of the boundary and cocycle exterior powers in a
free resolution.  The free degree-one term supplies
the line $Ou$.  Hence inside $Eu$ the full lattice is
\[
 \det_O^{-1}\mathscr C_\alpha
  =\pi^{\operatorname{length}H^2-
                    \operatorname{length}H^1_{\rm tors}}Ou.
                                                            \tag{EU.3}
\]
Since $z_\alpha$ generates that lattice, (EU.1) follows.
Changing $u$ multiplies $d_\alpha$ by a unit, so the
valuation is unchanged.  If $i_\alpha\geq0$, the
difference in (EU.2) equals $\operatorname{length}H^1_{\rm tors}$.
If $i_\alpha<0$, that difference is
$\operatorname{length}H^2=i_\alpha+
\operatorname{length}H^1_{\rm tors}$, which is
nonnegative and no larger than the torsion length.
\end{proof}

\begin{proposition}[AC-K26-elliptic: application to every fixed-base character]
\label{AC-K26-elliptic}
For every $\ell,\eta,e$ in AC-K25, put
\[
       T_{\ell,\eta,e}=O_\ell(\eta^{-1})(e+1),
       \quad
       C_{\ell,\eta,e}=
                 R\Gamma(G_{K,\Sigma_\ell},T_{\ell,\eta,e}).
\]
There is a specified, nonzero derived elliptic-unit
determinant element
$z_{\ell,\eta,e}\in H^1(C_{\ell,\eta,e})\otimes_{O_\ell}E_\ell$
to which AC-K26 applies.  Consequently, writing
$i_{\ell,\eta,e}$ for its index in the free
degree-one lattice,
\[
 \log\#H^2(C_{\ell,\eta,e})
  =\log\#k_{O_\ell}\,
       \bigl(i_{\ell,\eta,e}+
          \operatorname{length}_{O_\ell}H^1(C_{\ell,\eta,e})_{\rm tors}\bigr).
                                                            \tag{EU.4}
\]
This statement uses the characterwise main conjecture
also at $2$ and $3$.  It does not assume an equivariant
$\mu$-vanishing theorem there, and it asserts no
numerical upper bound for the indices.
\end{proposition}
\begin{proof}
Write the cyclotomic character as
$\kappa_\ell=\omega_\ell\langle\kappa_\ell\rangle$,
where $\omega_\ell$ has finite order and the second
factor has torsion-free pro-$\ell$ image.
For $\ell=2$ use the sign modulo $4$ and the factor
in $1+4\mathbb Z_2$.
The pro-$\ell$ factor cuts out the cyclotomic
$\mathbb Z_\ell$ extension of $K$, contained in its
$\mathbb Z_\ell^2$ extension $K_\infty$.
Set $\theta=\eta^{-1}\omega_\ell^e$.
This is a finite-order character, unramified outside
$\Sigma_\ell$, with values in $O_\ell$.

Use the ordinary scalar-action convention on
$\Lambda=O_\ell[[\Gamma]]$, where
$\Gamma=\operatorname{Gal}(K_\infty/K)$.
Take the universal rank-one representation whose
$G_K$ action is
\[
         g\longmapsto \theta(g)\kappa_\ell(g)[\bar g].
                                                            \tag{EU.5}
\]
Its specialization at
$\alpha_e([\gamma])=\langle\kappa_\ell\rangle(\gamma)^e$
is precisely $T_{\ell,\eta,e}$.
Transporting the source's scalar convention by inversion
is an automorphism of $\Lambda$; it preserves perfect
complexes, determinant isomorphisms and their base change.
The imprimitive characterwise theorem, with every prime
of $\Sigma_\ell$ removed, gives a specified generator
$z_\Lambda$ of the inverse determinant of the universal
cohomology complex.  In the source's convention this is
the elliptic-unit element of (17), multiplied by its
imprimitive Euler factors; Lemma~6.6 identifies the
determinant of its auxiliary ideal with the free line
generated by that element.
Apply its derived base-change assertion, with the
twisting automorphism if needed, to $\alpha_e$.
This gives exactly the displayed complex $C_{\ell,\eta,e}$.

We verify the finite-level hypotheses, rather than
infer them from a specialization of characteristic ideals.
The positive Tate weight and finite character imply
$H^0(C_{\ell,\eta,e})=0$: after killing the finite
character, the cyclotomic image is infinite and a
nonzero vector cannot be invariant under its
$(e+1)$st power.  There is no cohomology above degree
two for this totally imaginary arithmetic curve.
The degree-two rational cohomology vanishes.
Indeed, restriction to a finite field killing $\eta$
is injective rationally, by corestriction.  Over that
field the assertion at the primes above $\ell$ is
the finite even-K-group comparison (K10.5).
Adding primes away from $\ell$ changes degree two
by residue-field groups with Frobenius eigenvalue
$Nv^e$; these are torsion because $e\geq1$ and
$Nv^e\ne1$.  Thus $H^2$ is finite.

Finally the perfect universal complex has Euler
rank $-1$, by the rank-one and torsion assertions of
the characterwise theorem.  Removing additional primes
adds only two-term local complexes of Euler rank zero.
The alternating ranks of a finite free representative
do not change under specialization.  Since the only
remaining rational cohomology is in degree one,
it has rank one.  All hypotheses of AC-K26 hold.
Specialize the determinant generator itself and use
the rational cohomology identification to define
$z_{\ell,\eta,e}$.  Equations (EU.1) and (EU.4) follow.
This definition does not identify a naive tensor
product of $H^1$ with the specialized $H^1$, and
therefore remains valid when higher Tor terms occur.
\end{proof}

\begin{proposition}[AC-K26-size: a determinant identity does not bound its indices]
\label{AC-K26-size}
The rank, determinant, perfectness and $\mu=0$
assertions alone do not give a uniform upper bound
for the specialized lengths in (EU.1), even at a
fixed arithmetic point of a two-variable algebra.
\end{proposition}
\begin{proof}
Fix $a,b\in\pi O$, let $M\geq1$, and put
$f_M=T_1-a+\pi^M$ in $\Lambda=O[[T_1,T_2]]$.
Consider the perfect complex in degrees $1,2$
\[
 \Lambda^2\xrightarrow{(x,y)\mapsto f_Mx}\Lambda.
\]
Its degree-one cohomology is the free line
$\Lambda e_2$, its degree-two cohomology is
$\Lambda/(f_M)$, and its other cohomology vanishes.
The element $z_M=f_Me_2$ generates its inverse
determinant line.  Equivalently,
\[
 H^1/\Lambda z_M\simeq\Lambda/(f_M)\simeq H^2,
\]
so their characteristic divisors agree.
Since the coefficient of $T_1$ is one, $f_M$ is
not divisible by $\pi$; its $\mu$-invariant is zero.
Specialize at $T_1=a,T_2=b$.
The complex then has $H^1=Oe_2$, $H^2=O/\pi^MO$,
and $z_M=\pi^Me_2$.  Hence the index and the
degree-two length both equal the arbitrary integer $M$.
This is an algebraic model of precisely the listed
source conclusions, not an arithmetic counterexample
to a possible bound for elliptic units.
\end{proof}

\paragraph{Geometric meaning and remaining input.}
The universal Artin--Tate coefficient is a rank-one
sheaf on the fixed arithmetic curve over
$\mathcal O_K$ with a two-dimensional formal weight
space.  Its determinant is a line formed from a
perfect cohomology complex.  Equation (EU.1) separates
that line's integral specialization from its rational
rank.  This uses the determinant construction in the
cited paper's \S1.3; it is not a new claim that an
uninspected Grothendieck scan supplies a regulator bound.
The finite etale and K-theory dictionaries already
recorded for AC-K25 and AC-K10 still apply.
The remaining quantitative question is to bound the
sum of these actual elliptic-unit specialization
indices as the character, conductor, coefficient
prime and Tate weight vary together.  AC-K26-size
explains why this question cannot be omitted from
the main arithmetic argument.
